\section{Coherent sheaves on Noetherian schemes}
\label{section-coherent-quasi-compact}

\noindent
In this section we mention some properties of coherent sheaves on
Noetherian schemes.

\begin{lemma}
\label{lemma-acc-coherent}
Let $X$ be a Noetherian scheme.
Let $\mathcal{F}$ be a coherent $\mathcal{O}_X$-module.
The ascending chain condition holds for quasi-coherent submodules
of $\mathcal{F}$. In other words, given any sequence
$$
\mathcal{F}_1 \subset \mathcal{F}_2 \subset \ldots \subset \mathcal{F}
$$
of quasi-coherent submodules, then
$\mathcal{F}_n = \mathcal{F}_{n + 1} = \ldots $ for some $n \geq 0$.
\end{lemma}

\begin{proof}
Choose a finite affine open covering.
On each member of the covering we get stabilization by
Algebra, Lemma \ref{algebra-lemma-Noetherian-basic}.
Hence the lemma follows.
\end{proof}

\begin{lemma}
\label{lemma-power-ideal-kills-sheaf}
Let $X$ be a Noetherian scheme.
Let $\mathcal{F}$ be a coherent sheaf on $X$.
Let $\mathcal{I} \subset \mathcal{O}_X$ be a quasi-coherent
sheaf of ideals corresponding to a closed subscheme $Z \subset X$.
Then there is some $n \geq 0$ such that $\mathcal{I}^n\mathcal{F} = 0$
if and only if $\text{Supp}(\mathcal{F}) \subset Z$ (set theoretically).
\end{lemma}

\begin{proof}
This follows immediately from
Algebra, Lemma \ref{algebra-lemma-Noetherian-power-ideal-kills-module}
because $X$ has a finite covering by spectra of Noetherian rings.
\end{proof}

\begin{lemma}[Artin-Rees]
\label{lemma-Artin-Rees}
Let $X$ be a Noetherian scheme.
Let $\mathcal{F}$ be a coherent sheaf on $X$.
Let $\mathcal{G} \subset \mathcal{F}$ be a quasi-coherent subsheaf.
Let $\mathcal{I} \subset \mathcal{O}_X$ be a quasi-coherent sheaf of
ideals.
Then there exists a $c \geq 0$ such that for all $n \geq c$ we
have
$$
\mathcal{I}^{n - c}(\mathcal{I}^c\mathcal{F} \cap \mathcal{G})
=
\mathcal{I}^n\mathcal{F} \cap \mathcal{G}
$$
\end{lemma}

\begin{proof}
This follows immediately from
Algebra, Lemma \ref{algebra-lemma-Artin-Rees}
because $X$ has a finite covering by spectra of Noetherian rings.
\end{proof}

\begin{lemma}
\label{lemma-directed-colimit-coherent}
Let $X$ be a Noetherian scheme. Every quasi-coherent $\mathcal{O}_X$-module
is the filtered colimit of its coherent submodules.
\end{lemma}

\begin{proof}
This is a reformulation of Properties, Lemma
\ref{properties-lemma-quasi-coherent-colimit-finite-type}
in view of the fact that a finite type quasi-coherent
$\mathcal{O}_X$-module is coherent by
Lemma \ref{lemma-coherent-Noetherian}.
\end{proof}

\begin{lemma}
\label{lemma-homs-over-open}
Let $X$ be a Noetherian scheme.
Let $\mathcal{F}$ be a quasi-coherent $\mathcal{O}_X$-module.
Let $\mathcal{G}$ be a coherent $\mathcal{O}_X$-module.
Let $\mathcal{I} \subset \mathcal{O}_X$ be a quasi-coherent sheaf of
ideals. Denote $Z \subset X$ the corresponding closed subscheme and
set $U = X \setminus Z$.
There is a canonical isomorphism
$$
\colim_n \Hom_{\mathcal{O}_X}(\mathcal{I}^n\mathcal{G}, \mathcal{F})
\longrightarrow
\Hom_{\mathcal{O}_U}(\mathcal{G}|_U, \mathcal{F}|_U).
$$
In particular we have an isomorphism
$$
\colim_n \Hom_{\mathcal{O}_X}(
\mathcal{I}^n, \mathcal{F})
\longrightarrow
\Gamma(U, \mathcal{F}).
$$
\end{lemma}

\begin{proof}
We first prove the second map is an isomorphism. It is injective by 
Properties, Lemma \ref{properties-lemma-sections-over-quasi-compact-open}.
Since $\mathcal{F}$ is the union of its coherent submodules, see
Properties, Lemma \ref{properties-lemma-quasi-coherent-colimit-finite-type}
(and Lemma \ref{lemma-coherent-Noetherian})
we may and do assume that $\mathcal{F}$ is coherent to prove surjectivity.
Let $\mathcal{F}_n$ denote the quasi-coherent subsheaf of $\mathcal{F}$
consisting of sections annihilated by $\mathcal{I}^n$,
see Properties, Lemma \ref{properties-lemma-sections-over-quasi-compact-open}.
Since $\mathcal{F}_1 \subset \mathcal{F}_2 \subset \ldots$ we see that
$\mathcal{F}_n = \mathcal{F}_{n + 1} = \ldots $ for some $n \geq 0$
by Lemma \ref{lemma-acc-coherent}. Set $\mathcal{H} = \mathcal{F}_n$
for this $n$. By Artin-Rees (Lemma \ref{lemma-Artin-Rees})
there exists an $c \geq 0$ such that
$\mathcal{I}^m\mathcal{F} \cap \mathcal{H}
\subset \mathcal{I}^{m - c}\mathcal{H}$. Picking $m = n + c$ we get
$\mathcal{I}^m\mathcal{F} \cap \mathcal{H} \subset \mathcal{I}^n\mathcal{H}
= 0$. Thus if we set $\mathcal{F}' = \mathcal{I}^m\mathcal{F}$ then we
see that $\mathcal{F}' \cap \mathcal{F}_n = 0$ and
$\mathcal{F}'|_U = \mathcal{F}|_U$. Note in particular that the subsheaf
$(\mathcal{F}')_N$ of sections annihilated by $\mathcal{I}^N$ is zero
for all $N \geq 0$. Hence by
Properties, Lemma \ref{properties-lemma-sections-over-quasi-compact-open}
we deduce that
the top horizontal arrow in the following commutative
diagram is a bijection:
$$
\xymatrix{
\colim_n \Hom_{\mathcal{O}_X}(
\mathcal{I}^n, \mathcal{F}')
\ar[r] \ar[d] &
\Gamma(U, \mathcal{F}') \ar[d] \\
\colim_n \Hom_{\mathcal{O}_X}(
\mathcal{I}^n, \mathcal{F})
\ar[r] &
\Gamma(U, \mathcal{F})
}
$$
Since also the right vertical arrow is a bijection we conclude that
the bottom horizontal arrow is surjective as desired.

\medskip\noindent
Next, we prove the first arrow of the lemma is a bijection.
By Lemma \ref{lemma-coherent-Noetherian} the sheaf $\mathcal{G}$
is of finite presentation and hence the sheaf
$\mathcal{H} = \SheafHom_{\mathcal{O}_X}(\mathcal{G}, \mathcal{F})$
is quasi-coherent, see
Schemes, Section \ref{schemes-section-quasi-coherent}.
By definition we have
$$
\mathcal{H}(U)
=
\Hom_{\mathcal{O}_U}(\mathcal{G}|_U, \mathcal{F}|_U)
$$
Pick a $\psi$ in the right hand side of the first arrow of the
lemma, i.e.,  $\psi \in \mathcal{H}(U)$. The result just proved applies
to $\mathcal{H}$ and hence there exists an $n \geq 0$ and an
$\varphi : \mathcal{I}^n \to \mathcal{H}$ which recovers
$\psi$ on restriction to $U$. By
Modules, Lemma \ref{modules-lemma-internal-hom}
$\varphi$ corresponds to a map
$$
\varphi :
\mathcal{I}^n \otimes_{\mathcal{O}_X} \mathcal{G}
\longrightarrow
\mathcal{F}.
$$
This is almost what we want except that the source of the arrow
is the tensor product of $\mathcal{I}^n$ and $\mathcal{G}$
and not the product. We will show that, at the cost of increasing $n$,
the difference is irrelevant. Consider the short exact sequence
$$
0 \to \mathcal{K} \to
\mathcal{I}^n \otimes_{\mathcal{O}_X} \mathcal{G} \to
\mathcal{I}^n\mathcal{G} \to 0
$$
where $\mathcal{K}$ is defined as the kernel. Note that
$\mathcal{I}^n\mathcal{K} = 0$ (proof omitted). By Artin-Rees
again we see that
$$
\mathcal{K}
\cap
\mathcal{I}^m(\mathcal{I}^n \otimes_{\mathcal{O}_X} \mathcal{G})
=
0
$$
for some $m$ large enough. In other words we see that
$$
\mathcal{I}^m(\mathcal{I}^n \otimes_{\mathcal{O}_X} \mathcal{G})
\longrightarrow
\mathcal{I}^{n + m}\mathcal{G}
$$
is an isomorphism. Let $\varphi'$ be the restriction of
$\varphi$ to this submodule thought of as a map
$\mathcal{I}^{m + n}\mathcal{G} \to \mathcal{F}$.
Then $\varphi'$ gives an element
of the left hand side of the first arrow of the lemma which
maps to $\psi$ via the arrow. In other words we have proved surjectivity
of the arrow. We omit the proof of injectivity.
\end{proof}

\begin{lemma}
\label{lemma-extend-coherent}
Let $X$ be a locally Noetherian scheme. Let $\mathcal{F}$, $\mathcal{G}$
be coherent $\mathcal{O}_X$-modules. Let $U \subset X$ be open and
let $\varphi : \mathcal{F}|_U \to \mathcal{G}|_U$ be an
$\mathcal{O}_U$-module map. Then there exists a coherent
submodule $\mathcal{F}' \subset \mathcal{F}$ agreeing with
$\mathcal{F}$ over $U$ such that $\varphi$ extends to
$\varphi' : \mathcal{F}' \to \mathcal{G}$.
\end{lemma}

\begin{proof}
Let $\mathcal{I} \subset \mathcal{O}_X$ be the coherent sheaf of ideals
cutting out the reduced induced scheme structure on $X \setminus U$.
If $X$ is Noetherian, then Lemma \ref{lemma-homs-over-open} tells us that
we can take $\mathcal{F}' = \mathcal{I}^n\mathcal{F}$
for some $n$. The general case will follow from this using Zorn's lemma.

\medskip\noindent
Consider the set of triples $(U', \mathcal{F}', \varphi')$ where
$U \subset U' \subset X$ is open, $\mathcal{F}' \subset \mathcal{F}|_{U'}$
is a coherent subsheaf agreeing with $\mathcal{F}$ over $U$, and
$\varphi' : \mathcal{F}' \to \mathcal{G}|_{U'}$ restricts to
$\varphi$ over $U$. We say
$(U'', \mathcal{F}'', \varphi'') \geq (U', \mathcal{F}', \varphi')$
if and only if $U'' \supset U'$, $\mathcal{F}''|_{U'} = \mathcal{F}'$,
and $\varphi''|_{U'} = \varphi'$.
It is clear that if we have a totally ordered collection
of triples $(U_i, \mathcal{F}_i, \varphi_i)$, then we
can glue the $\mathcal{F}_i$ to a subsheaf $\mathcal{F}'$ of $\mathcal{F}$
over $U' = \bigcup U_i$ and extend $\varphi$ to a map
$\varphi' : \mathcal{F}' \to \mathcal{G}|_{U'}$.
Hence any totally ordered subset of triples has an upper bound.
Finally, suppose that $(U', \mathcal{F}', \varphi')$
is any triple but $U' \not = X$. Then we can choose an
affine open $W \subset X$ which is not contained in $U'$.
By the result of the first paragraph we can extend
the subsheaf $\mathcal{F}'|_{W \cap U'}$ and the restriction
$\varphi'|_{W \cap U'}$ to some subsheaf $\mathcal{F}'' \subset \mathcal{F}|_W$
and map $\varphi'' : \mathcal{F}'' \to \mathcal{G}|_W$.
Of course the agreement between $(\mathcal{F}', \varphi')$ and
$(\mathcal{F}'', \varphi'')$ over $W \cap U'$ exactly means that
we can extend this to a triple $(U' \cup W, \mathcal{F}''', \varphi''')$.
Hence any maximal triple $(U', \mathcal{F}', \varphi')$
(which exist by Zorn's lemma) must have $U' = X$ and the
proof is complete.
\end{proof}









